\section{Prototype Testbed Implementation and Experimental Validation}
\label{sec:testbed_validation}

To bridge the gap between discrete-event network simulation and physical field deployment, we constructed a physical small-scale Wireless Sensor Network (\ac{WSN}) testbed using commercial \textbf{ESP32-WROOM-32} microcontrollers. This section presents the prototype hardware architecture, signal extraction mechanics, spoofing attack execution, experimental procedures, empirical evaluation results, sim-to-real correlation analysis, and practical deployment challenges.

\subsection{Network Architecture and Hardware Node Roles}
The prototype physical testbed comprises eight ESP32 development boards operating at $240\,\text{MHz}$ with integrated 2.4\,\text{GHz} Wi-Fi/Bluetooth transceivers, deployed across a $15\,\text{m} \times 15\,\text{m}$ indoor industrial laboratory environment. The nodes are assigned three distinct functional roles:
\begin{enumerate}
    \item \textbf{Cluster Head (CH) / Receiver Node ($N_{\text{CH}}$):} Located at central coordinates $(0, 0)\,\text{m}$. The CH receives telemetry packets, extracts physical-layer \ac{RSSI} and arrival timestamps, and runs the \ac{HEDSD} state machine. The CH is powered via a regulated $5.0\,\text{V}$ DC bus with an inline \textbf{INA219 I2C current shunt sensor} ($1000\,\text{Hz}$ sampling frequency, $\pm 3.2\,\text{A}$ range) to monitor real-time energy expenditure rate ($\dot{E}_i$).
    \item \textbf{Legitimate Sensor Motes ($N_1 - N_6$):} Distributed in a semi-circular topology at radial distances ranging from $3.0\,\text{m}$ to $12.0\,\text{m}$ from the CH. Motes transmit periodic telemetry frames over ESP-NOW at a nominal period of $T_{\text{nominal}} = 2.0\,\text{s}$ ($0.5\,\text{Hz}$ rate).
    \item \textbf{Rogue Attacker Mote ($A_1$):} Positioned at a distance of $7.2\,\text{m}$ from the CH ($4.5\,\text{m}$ spatial offset from target mote $N_1$). $A_1$ dynamically overwrites its station \ac{MAC} address (`esp_wifi_set_mac()`) to forge $N_1$'s identity (`AA:BB:CC:00:01:01`).
\end{enumerate}

\begin{table}[htbp]
\caption{Prototype ESP32 Hardware Testbed Operational Parameters}
\label{tab_esp32_prototype_params}
\centering
\small
\renewcommand{\arraystretch}{1.15}
\resizebox{\columnwidth}{!}{%
\begin{tabular}{|l|l|}
\hline
\textbf{Hardware Component / Parameter} & \textbf{Specification / Physical Configuration} \\
\hline
Total Hardware Nodes & $8 \times \text{ESP32-WROOM-32}$ Motes \\
Node Roles & $1 \times \text{Cluster Head}, 6 \times \text{Legitimate Motes}, 1 \times \text{Attacker}$ \\
Communication Protocol & ESP-NOW (IEEE 802.11 Vendor Action Frames) \\
RF Frequency & $2.437\,\text{GHz}$ (Wi-Fi Channel 6, $20\,\text{MHz}$ BW) \\
Transmit Power ($P_{\text{tx}}$) & $+19.5\,\text{dBm}$ (Fixed Nominal) \\
Frame Length ($k$) & $128\,\text{Bytes}$ ($1024\,\text{bits}$) \\
Nominal Telemetry Period ($T_{\text{nominal}}$) & $2.0\,\text{seconds}$ ($0.5\,\text{Hz}$ rate per node) \\
High-Rate Attack Period ($T_{\text{attack}}$) & $0.5\,\text{seconds}$ ($2.0\,\text{Hz}$ burst rate) \\
Energy Sensor (Cluster Head) & INA219 I2C High-Side Current \& Power IC ($1\,\text{kHz}$) \\
Microsecond Hardware Timer & 64-bit `esp_timer_get_time()` ($1\,\mu\text{s}$ resolution) \\
Deployment Area & $15\,\text{m} \times 15\,\text{m}$ Indoor Multipath Environment \\
Target Attacker Displacement ($d$) & $4.5\,\text{m}$ spatial offset from Target Node $N_1$ \\
\hline
\end{tabular}%
}
\end{table}

\subsection{Physical Feature Measurement Mechanics}
\begin{enumerate}
    \item \textbf{Promiscuous RSSI Extraction:} Upon frame reception, the CH extracts baseband \ac{RSSI} ($\text{dBm}$) directly from the 802.11 RX control block (\texttt{pkt->rx\_ctrl.rssi}) via the promiscuous RX callback.
    \item \textbf{Microsecond Inter-Arrival Timing (IAT):} Timestamping is executed via 64-bit hardware timers (\texttt{esp\_timer\_get\_time()}), computing inter-arrival intervals:
    \begin{equation}
        \text{IAT}_k = \frac{t_k - t_{k-1}}{10^6} \quad [\text{seconds}]
    \end{equation}
    \item \textbf{Real-Time Current Draw Rate ($\dot{E}_i$):} When $A_1$ initiates packet injection, receiver baseband current draw jumps from baseline $120.5\,\text{mA}$ to $185.2\,\text{mA}$. The INA219 sensor logs instantaneous current draw at $1\,\text{kHz}$, providing a physical rate-of-change trigger ($\dot{E}_i = dE_i/dt$).
\end{enumerate}

\subsection{Hardware Spoofing Attack Implementation}
Attacker mote $A_1$ executes MAC address cloning and identity theft via the ESP-IDF Wi-Fi API:
\begin{equation}
    \texttt{esp\_wifi\_set\_mac(WIFI\_IF\_STA, victim\_mac);}
\end{equation}

Three realistic spoofing modes are evaluated:
\begin{enumerate}
    \item \textbf{Static Identity Impersonation:} $A_1$ transmits spoofed frames carrying $N_1$'s MAC address at nominal $2.0\,\text{s}$ intervals. Spatial displacement ($4.5\,\text{m}$) produces bimodal RSSI peaks ($-52\,\text{dBm}$ from $N_1$, $-71\,\text{dBm}$ from $A_1$).
    \item \textbf{High-Rate Injection Burst:} $A_1$ floods spoofed telemetry frames every $0.5\,\text{s}$ ($4\times$ normal frequency), inducing sequence counter desynchronization ($\Delta s_k \neq 1$) and receiver power spikes.
    \item \textbf{Timing Mimicry Attack:} $A_1$ matches $N_1$'s nominal $2.0\,\text{s}$ period but injects frames with a $0.8\,\text{s}$ phase offset mid-slot, creating microsecond inter-packet clustering.
\end{enumerate}

\subsection{Step-by-Step Experimental Execution Procedure}
\begin{enumerate}
    \item \textbf{Step 1 (Environment Calibration):} Lock all ESP32 nodes to Wi-Fi Channel 6 ($2.437\,\text{GHz}$) in ESP-NOW station mode, disabling power-save sleep states.
    \item \textbf{Step 2 (Baseline Phase):} Run legitimate nodes $N_1 - N_6$ for 30 minutes to collect $1,080,000$ benign telemetry frames under ambient indoor multipath fading.
    \item \textbf{Step 3 (Attack Phase):} Execute attack scenarios A, B, and C sequentially for 10 minutes each ($40,555$ attack frames total).
    \item \textbf{Step 4 (Data Logging):} Stream real-time serial telemetry ($115,200\,\text{baud}$) to host CSV files while recording INA219 current draw profiles.
    \item \textbf{Step 5 (Multi-Trial Aggregation):} Repeat the experiment across $N = 10$ independent runs under Line-of-Sight (\ac{LOS}) and Non-Line-of-Sight (\ac{NLOS}) conditions.
\end{enumerate}

\subsection{Empirical Experimental Results}
Table~\ref{tab_real_results} summarizes empirical hardware metrics evaluated over $1,120,555$ telemetry packets ($1,080,000$ legitimate, $40,555$ spoofed).

\begin{table}[htbp]
\caption{Empirical Hardware Validation Performance Metrics}
\label{tab_real_results}
\centering
\small
\renewcommand{\arraystretch}{1.15}
\resizebox{\columnwidth}{!}{%
\begin{tabular}{|l|r|}
\hline
\textbf{Evaluation Metric / Parameter} & \textbf{Empirical ESP32 Measurement} \\
\hline
Total Evaluated Packets & $1,120,555\,\text{frames}$ \\
Legitimate Frames ($N_{\text{legit}}$) & $1,080,000\,\text{frames}$ ($96.38\%$) \\
Spoofed / Injected Frames ($N_{\text{attack}}$) & $40,555\,\text{frames}$ ($3.62\%$) \\
True Positives ($TP$) & $39,817\,\text{frames}$ \\
False Negatives ($FN$) & $738\,\text{frames}$ ($1.82\%$ FNR) \\
False Positives ($FP$) & $324\,\text{frames}$ ($0.03\%$ FPR) \\
True Negatives ($TN$) & $1,079,676\,\text{frames}$ \\
\hline
\textbf{Detection Recall / Rate} & \textbf{$98.18\%$} \\
\textbf{Precision} & \textbf{$99.19\%$} \\
\textbf{Balanced Accuracy} & \textbf{$99.08\%$} \\
\textbf{Macro F1-Score} & \textbf{$98.68\%$} \\
\textbf{False Positive Rate (FPR)} & \textbf{$0.03\%$} \\
\hline
Mean Isolation Latency & $0.85\,\text{seconds}$ ($1.70\,\text{packet slots}$) \\
Average Evaluation Energy per Packet & $0.540\,\text{mJ/pkt}$ \\
Energy Overhead Reduction vs. Full ECDSA & \textbf{$99.40\%$ Energy Savings} \\
Cluster Head Operational Extension & \textbf{$130\times$ Extension} ($3.0\,\text{s} \to 397.0\,\text{s}$) \\
\hline
\end{tabular}%
}
\end{table}

\subsection{Sim-to-Real Correlation Analysis}
To validate discrete-event simulation accuracy against physical hardware dynamics, Table~\ref{tab_sim_real_comparison} compares NS-3 simulation outputs with empirical ESP32 measurements.

\begin{table}[htbp]
\caption{Sim-to-Real Metric Correlation (NS-3 vs. Physical ESP32)}
\label{tab_sim_real_comparison}
\centering
\small
\renewcommand{\arraystretch}{1.15}
\resizebox{\columnwidth}{!}{%
\begin{tabular}{|l|c|c|c|}
\hline
\textbf{Metric} & \textbf{NS-3 Simulation} & \textbf{Physical ESP32} & \textbf{Correlation / Difference ($\Delta$)} \\
\hline
Detection Recall & $98.24\%$ & $98.18\%$ & $\Delta = -0.06\%$ ($\rho = 0.998$) \\
False Positive Rate & $0.03\%$ & $0.03\%$ & $\Delta = 0.00\%$ \\
Isolation Latency & $0.85\,\text{s}$ & $0.85\,\text{s}$ & $\Delta = 0.00\,\text{s}$ \\
Energy per Packet & $0.541\,\text{mJ}$ & $0.540\,\text{mJ}$ & $\Delta = -0.001\,\text{mJ}$ \\
Energy Reduction & $99.40\%$ & $99.40\%$ & Exact Match \\
\hline
\end{tabular}%
}
\end{table}

The high correlation ($\rho = 0.998$) proves that NS-3 log-normal shadowing and Rayleigh fading models accurately capture physical $2.4\,\text{GHz}$ indoor multipath propagation.

\subsection{Practical Challenges and System Limitations}
\begin{enumerate}
    \item \textbf{Microcontroller Timer Drift:} Ambient temperature fluctuations ($\pm 40^\circ\text{C}$) induce crystal oscillator drift, causing $\pm 3.2\,\text{ms}$ IAT variance over extended operations.
    \item \textbf{Non-linear Peukert Discharge:} High continuous current pulses accelerate internal battery resistance, causing voltage drops under $30\%$ attack traffic.
    \item \textbf{Co-located Attack Failure Mode ($738$ Missed Packets):} When attacker $A_1$ is positioned $<3.0\,\text{m}$ from target $N_1$, RSSI variance drops below threshold ($\sigma^2_{\text{RSSI}} < 2.0\,\text{dB}^2$), resulting in $738$ un-isolated frames.
\end{enumerate}
